\documentclass[11pt,a4paper]{article}
\usepackage[margin=20mm]{geometry}
\usepackage{amsmath,amssymb,amsthm,mathtools,enumitem,xcolor,tcolorbox,hyperref,fancyhdr,microtype,longtable}
\tcbuselibrary{breakable,skins}
\definecolor{navy}{HTML}{15304D}\definecolor{sky}{HTML}{EAF4FB}\definecolor{olive}{HTML}{17645F}
\hypersetup{colorlinks=true,linkcolor=navy,urlcolor=olive}
\pagestyle{fancy}\fancyhf{}\lhead{MATH6310 | Metric Spaces \& Topology}\rhead{Sections 0--1}\cfoot{\thepage}
\setlength{\parskip}{5pt}\setlength{\parindent}{0pt}
\newtcolorbox{keybox}[1]{enhanced,breakable,colback=sky,colframe=navy,title=\textbf{#1},fonttitle=\color{white},colbacktitle=navy}
\newtcolorbox{tipbox}[1]{enhanced,breakable,colback=white,colframe=olive,title=\textbf{#1},fonttitle=\color{white},colbacktitle=olive}
\newcommand{\R}{\mathbb R}\newcommand{\N}{\mathbb N}\newcommand{\Z}{\mathbb Z}\newcommand{\Q}{\mathbb Q}\newcommand{\Pow}{\mathcal P}\newcommand{\dd}{\mathrm d}\newcommand{\im}{\operatorname{Ran}}\newcommand{\id}{\operatorname{id}}
\newcommand{\ex}[2]{\subsection*{Exercise #1: #2}\addcontentsline{toc}{subsection}{Exercise #1: #2}}
\begin{document}
\begin{titlepage}\centering\vspace*{25mm}{\Huge\bfseries Metric Spaces \& Topology\par}\vspace{5mm}{\LARGE MATH6310\par}\vspace{10mm}{\Large Expanded revision textbook, intuitive explanations and mark-by-mark worked solutions\par}\vspace{7mm}{\large Sections 0 and 1: Metrics, sets, functions, equivalence relations, countability and orders\par}\vfill\begin{tcolorbox}[colback=sky,colframe=navy,width=.88\textwidth]\centering\large\textbf{Understand the idea $\to$ learn the definition $\to$ recognise the proof pattern $\to$ solve the question.}\end{tcolorbox}\vfill{\large Based on both MATH6310 lecture-note and exercise PDFs dated 7 September 2026.}\par\vspace{8mm}\textit{A student-friendly companion; proofs and worked answers are newly written, not transcriptions of the empty solution slides.}\end{titlepage}
\tableofcontents\newpage
\section{How to use this guide}
This expanded edition gives plain-English explanations, exam-writing tips, and suggested stepwise mark allocations. Where a mark total is printed in the source notes it is retained; otherwise allocations are explicitly illustrative, not official. This guide covers the mathematical definitions, examples and main theorems in the two supplied documents, then gives worked solutions to \textbf{all six numbered Section 0 exercises and all nineteen numbered Section 1 exercises}. Exercises repeated inside the lecture notes are answered in their numbered exercise counterparts; two extra image/difference identities appearing in the teaching slides are proved with Exercise 1.8. An optional-theory review is included for Schr\"oder--Bernstein and Zorn's lemma.
\begin{tipbox}{A good proof-exam strategy}Write down the definition \emph{before} calculating. For a metric, check non-negativity/finiteness, identity of indiscernibles, symmetry, and the triangle inequality. For an equality of sets, take an arbitrary element and show membership in both directions. For a bijection, give an explicit inverse and verify both compositions.\end{tipbox}
\section{Section 0: The metric-space cheat sheet}
\subsection{What is the point of metric spaces?}
In $\R$, distance usually means $|x-y|$. In $\R^n$, it might be straight-line distance. For functions, we may measure the biggest difference between two graphs, or the total area between them. A \emph{metric space} packages all these possibilities into one language. The same proofs of continuity and convergence can then be used in many settings.
\begin{keybox}{Metric: learn this precisely}A metric on a set $S$ is $d:S\times S\to[0,\infty)$ with, for all $x,y,z\in S$,
\begin{align*} \text{(MS1)}&\quad d(x,y)=0\iff x=y,\\ \text{(MS2)}&\quad d(x,y)=d(y,x),\\ \text{(MS3)}&\quad d(x,z)\le d(x,y)+d(y,z).\end{align*}
The codomain $[0,\infty)$ supplies \textbf{non-negativity and finiteness}. The notes also split MS1 into \textbf{MS1a} $d(x,x)=0$ and \textbf{MS1b} $d(x,y)=0\Rightarrow x=y$. A metric space is the pair $(S,d)$, not merely the set $S$.\end{keybox}
\subsection{The important metric examples}
\begin{description}[leftmargin=18mm,style=nextline]
\item[Standard metric on $\R$] $d(x,y)=|x-y|$. Its triangle inequality is $|a+b|\le|a|+|b|$.
\item[Euclidean metric on $\R^n$] $d_2(x,y)=\|x-y\|_2=\sqrt{\sum_{i=1}^n(x_i-y_i)^2}$. The norm triangle inequality follows from Cauchy--Schwarz.
\item[Discrete metric on any set] $d_{\rm disc}(x,y)=0$ if $x=y$, and $1$ otherwise. It ignores how \emph{far} distinct points look geometrically.
\item[Counting metric on $\Pow(S)$, $S$ finite] $d(A,B)=|A\triangle B|$, where $A\triangle B=(A\setminus B)\cup(B\setminus A)$. It counts membership disagreements.
\item[Hausdorff metric] On nonempty bounded closed sets $K,J\subset\R^n$, define $K^\delta=\{x:\exists k\in K,\ \|x-k\|_2\le\delta\}$ and $d_H(K,J)=\inf\{\delta\ge0:K^\delta\supseteq J,\ J^\delta\supseteq K\}$. Intuition: how much must each set be thickened to cover the other?
\item[Integration metric on continuous functions] On $C([0,1];\R)$, $d_1(f,g)=\int_0^1|f(t)-g(t)|\,dt$. Measures total absolute discrepancy. Continuity is essential for positive definiteness.
\item[Uniform/supremum metric] $d_\infty(f,g)=\sup_{t\in[0,1]}|f(t)-g(t)|$. Measures worst-case discrepancy. Continuity on the compact interval ensures finiteness.
\item[Two unusual metrics on $\R^2$] $d_{\rm TGV}(u,v)=0$ if $u=v$ and $\|u\|+\|v\|$ otherwise. The variant $d'_{\rm TGV}(u,v)=\|u-v\|$ when $u,v$ are linearly dependent, and $\|u\|+\|v\|$ when independent.
\end{description}
\begin{tipbox}{Think before using the ordinary triangle inequality}If a distance is defined in cases, split your proof into cases. If it is defined by an integral, use $|f-h|\le|f-g|+|g-h|$ pointwise \emph{then integrate}. If it is defined by a supremum, use the same inequality \emph{then take suprema}. If it counts differing positions or set membership, show that every difference between the endpoints must appear in at least one intermediate difference.\end{tipbox}
\subsection{Useful worked examples}
\textbf{Example: counting metric.} For $S=\{1,2,3,4\}$, $A=\{1,3\}$ and $B=\{2,3\}$, the symmetric difference is $\{1,2\}$: the element $3$ is in both, so it is not counted. Hence $d(A,B)=2$.

\textbf{Example: integration versus sup metric.} Let $f(t)=0$ and $g(t)=t$ on $[0,1]$. Then $d_1(f,g)=\int_0^1t\,dt=\tfrac12$, whereas $d_\infty(f,g)=\sup_{[0,1]}t=1$. Neither formula is a mistake: they measure different things.
\section{Section 1: Sets and functions cheat sheet}
\subsection{Sets and logic}
\begin{keybox}{Set-notation essentials}
$A\subseteq B$ means $\forall x\;(x\in A\Rightarrow x\in B)$; $A\subsetneq B$ means additionally $A\neq B$; $A\setminus B=\{x\in A:x\notin B\}$; $\Pow(S)=\{A:A\subseteq S\}$. The \textbf{symmetric difference} $A\triangle B=(A\setminus B)\cup(B\setminus A)$ contains the elements belonging to exactly one of $A,B$.\end{keybox}
For a family $\mathcal C$ of subsets of a universe $S$,
\[x\in\bigcup\mathcal C\iff(\exists C\in\mathcal C)\ x\in C,\qquad x\in\bigcap\mathcal C\iff(\forall C\in\mathcal C)\ x\in C.\]
For an empty family, $\bigcup\varnothing=\varnothing$ and $\bigcap\varnothing=S$ (relative to the stated universe). Important identities:
\begin{align*}A\cap\bigcup_{C\in\mathcal C}C&=\bigcup_{C\in\mathcal C}(A\cap C),\\ A\setminus\bigcup_{C\in\mathcal C}C&=\bigcap_{C\in\mathcal C}(A\setminus C),\\ A\setminus\bigcap_{C\in\mathcal C}C&=\bigcup_{C\in\mathcal C}(A\setminus C).\end{align*}
\subsection{Products and the axiom of choice}
$S\times T=\{(s,t):s\in S,t\in T\}$. An indexed Cartesian product $\prod_{\lambda\in\Lambda}S_\lambda$ consists of all families $(s_\lambda)$ such that $s_\lambda\in S_\lambda$ for every $\lambda$. It has one empty family if $\Lambda=\varnothing$; it is empty if some factor is empty. \textbf{Axiom of choice (AC):} any indexed family of nonempty sets admits a choice function $c(\lambda)\in S_\lambda$. In a substantial infinite family, choosing one element from each set may require AC.
\subsection{Functions, domain, codomain, range}
A function $f:S\to T$ specifies a value $f(s)\in T$ for every $s\in S$. Its \emph{range} is $\im f=\{f(s):s\in S\}\subseteq T$, which can be smaller than its codomain. The notes write the set of maps as $\mathcal F(S;T)$.

Useful examples: $\id_S(s)=s$; inclusion $\iota:A\hookrightarrow S$ sends $a\mapsto a$; indicator $1_A(s)=1$ on $A$ and $0$ otherwise; a constant map sends everything to one fixed point. For finite $|S|=n, |T|=m$, the number of maps $S\to T$ is $m^n$. An empty domain has exactly one function into every $T$; a nonempty domain has no function into $\varnothing$.
\begin{keybox}{Injective, surjective, bijective}
\textbf{Injective:} $f(x)=f(y)\Rightarrow x=y$ (different inputs cannot collide). \textbf{Surjective:} every $t\in T$ has a preimage in $S$ (range $=T$). \textbf{Bijective:} both, so every output has exactly one input. \textbf{Composition:} $(g\circ f)(s)=g(f(s))$: apply the right-hand function first.\end{keybox}
A restriction $f|_A:A\to T$ only changes the domain; a \emph{co-restriction} $f:S\to B$, $\im f\subseteq B\subseteq T$, changes the codomain. Inverses: a left inverse $g\circ f=\id_S$ implies injectivity; a right inverse $f\circ h=\id_T$ implies surjectivity. Conversely, for nonempty domain the left inverse can be constructed for an injection; a surjection has a right inverse using choice. A bijection has a unique two-sided inverse.
\subsection{Direct and inverse images}
For $A\subseteq S$, $f[A]=\{f(s):s\in A\}$ (direct image). For $C\subseteq T$, $f^{-1}[C]=\{s\in S:f(s)\in C\}$ (inverse image): \textbf{this inverse-image notation does not require an inverse function.}
\begin{align*}f[\bigcup_i A_i]&=\bigcup_i f[A_i],& f[\bigcap_i A_i]&\subseteq\bigcap_i f[A_i],\\ f[A\setminus B]&\supseteq f[A]\setminus f[B],& f^{-1}[\bigcup_i C_i]&=\bigcup_i f^{-1}[C_i],\\ f^{-1}[\bigcap_i C_i]&=\bigcap_i f^{-1}[C_i],& f^{-1}[D\setminus C]&=f^{-1}[D]\setminus f^{-1}[C].\end{align*}
The direct-image containments can be strict: if $f(0)=f(1)$ and $A=\{0\}$, $B=\{1\}$ then $f[A\cap B]=\varnothing$ while $f[A]\cap f[B]\ne\varnothing$.
\subsection{Equivalence relations and partitions}
An equivalence relation $\sim$ on $S$ is \textbf{reflexive} ($x\sim x$), \textbf{symmetric} ($x\sim y\Rightarrow y\sim x$), and \textbf{transitive} ($x\sim y\land y\sim z\Rightarrow x\sim z$). Its equivalence class is $[x]=\{y:y\sim x\}$; the quotient $S/{\sim}$ is the set of classes. Equivalence classes are either identical or disjoint and partition $S$. Conversely, every partition into nonempty disjoint blocks gives an equivalence relation: $x\sim y$ iff they lie in the same block. Example: $m\sim n$ when $p$ divides $m-n$.

\textbf{Canonical factorisation of a function.} The relation $x\sim y$ iff $f(x)=f(y)$ yields a quotient $S/{\sim}$; the map $S\to S/{\sim}$ is surjective, the induced map $S/{\sim}\to\im f$ is bijective, and $\im f\hookrightarrow T$ is injective. Their composition is $f$. This separates the many-to-one, one-to-one, and codomain features of any map.
\subsection{Countability and Cantor}
In the notes, a set is \emph{countable} if it is finite or can be listed by natural numbers; equivalently, it is empty or is the image of a surjection $\N\to S$. Basic facts: subsets of countable sets are countable; finite products of countable sets are countable; countable unions of countable sets are countable (the general statement uses choice). $\Z$ and $\Q$ are countable; a countable infinite set is bijective to $\N$.
\begin{keybox}{Cantor's theorem: a major result}For every set $S$, there is no surjection $S\to\Pow(S)$. The proof assumes a proposed list $f:S\to\Pow(S)$ and forms $D=\{s\in S:s\notin f(s)\}$. If $D=f(a)$, then $a\in D\iff a\notin D$, a contradiction. Consequently $\Pow(\N)$ and the set $\{0,1\}^{\N}$ are uncountable.\end{keybox}
\subsection{Orders and optional material}
A \emph{partial order} $\le$ is reflexive, antisymmetric and transitive. A \emph{total order} additionally compares any two elements. A strict total order $<$ is transitive and satisfies trichotomy: exactly one of $x<y$, $x=y$, $y<x$. A \emph{chain} is a totally ordered subset; a \emph{maximal} element has no strictly larger element in the poset.

\textbf{Zorn's lemma (optional section):} If a nonempty partially ordered set has an upper bound for every chain, it has a maximal element. Under the usual set-theoretic framework, it is equivalent to AC; it can be used to prove that every vector space has a Hamel basis.

\textbf{Schr\"oder--Bernstein:} If there are injections $S\to T$ and $T\to S$, there is a bijection $S\to T$. The notes explicitly require knowing the \emph{statement}, but say the proof is not examinable. Its key lesson is that two injections suffice to prove equality of cardinality.

\newpage
\section{First-principles teaching: understanding the concepts}
\subsection{How to read mathematical statements}
A definition is an instruction for checking whether an object has a property. Read $\forall$ as ``for every'' and $\exists$ as ``there exists''. For example, the metric triangle inequality means: \emph{no matter which three points you give me}, travelling from the first to the third directly is no longer than travelling via the second. A single counterexample is enough to disprove a universal statement.

The implication $P\Rightarrow Q$ does not normally mean $Q\Rightarrow P$. To prove $P\iff Q$, prove both directions. To prove two sets equal, show both inclusions or write a reversible ``if and only if'' membership chain. A contrapositive proof proves $\neg Q\Rightarrow\neg P$, logically equivalent to $P\Rightarrow Q$. A contradiction proof assumes the negation of the target and derives an impossibility.
\begin{tipbox}{How to earn proof marks}Give the general arbitrary object first, state the definition you are applying, justify the non-obvious step (often triangle inequality, continuity, or an element witness), and explicitly state the requested conclusion. An isolated final formula often cannot earn the reasoning marks.\end{tipbox}
\subsection{Metric axioms with small examples}
Let $S=\{a,b,c\}$. The discrete metric assigns distance $0$ from a point to itself and distance $1$ between all distinct points. For $a\ne c$, the triangle inequality reads $1\le d(a,b)+d(b,c)$: at least one leg is $1$. By contrast, the tempting formula $d(a,b)=2$ for distinct points except $d(a,c)=5$ fails if $a,b,c$ are distinct and $d(a,b)=d(b,c)=2$, because $5\nleq 2+2$.

For a metric proof, \textbf{MS0} verifies a finite nonnegative real value; \textbf{MS1a} verifies $d(x,x)=0$; \textbf{MS1b} verifies that no two distinct points have zero distance; \textbf{MS2} reverses the arguments; and \textbf{MS3} compares the direct and two-leg distances. The first three often take a line each; the triangle inequality generally needs the most explanation.
\subsection{Distances between functions are not all the same}
Suppose $f(t)=0$ and $g(t)=t^2$. The integral metric measures \emph{total accumulated difference}, $d_1(f,g)=\int_0^1t^2\,dt=1/3$. The sup metric measures the \emph{largest single difference}, $d_\infty(f,g)=1$. A tall, narrow discrepancy can have a small integral but a large supremum.

For the integral metric's definiteness, continuity is the central idea: if the function is nonzero even at one point, it must remain noticeably nonzero throughout a short nearby interval, creating strictly positive area. For the sup metric, zero supremum forces every single nonnegative difference to equal zero, and continuity is needed for boundedness, not for that last argument.
\subsection{What Hausdorff distance actually measures}
Imagine two shapes drawn on paper. Thicken each shape by a radius $r$, as though tracing with a pen of thickness $r$. The Hausdorff distance is the smallest $r$ allowing each thickened shape to cover the other original shape. This is a distance between \emph{sets}, not simply between their centres. The proof technique is a two-stage journey: a point of set $L$ is within $\delta$ of a point in $J$, which is within $\gamma$ of a point in $K$; therefore it is within $\gamma+\delta$ of $K$.

Take singletons $K=\{(0,0)\}$ and $J=\{(3,4)\}$. Their Hausdorff distance is $5$, precisely their ordinary Euclidean distance. This sanity check links the abstract formula back to geometry.
\subsection{Symmetric difference as an error detector}
$A\triangle B$ lists elements on which the answers to ``does this belong to the set?'' disagree. On $S=\{1,2,3,4\}$, $A=\{1,3\}$ and $B=\{2,3\}$ disagree about $1$ and $2$, so their counting distance is $2$. The triangle inequality follows because any membership disagreement between $A$ and $C$ must occur between $A$ and $B$ or between $B$ and $C$ (possibly both).
\subsection{Domain, codomain and range: three different jobs}
Consider $f:\R\to\R$, $f(x)=x^2$. The \emph{domain} tells you valid inputs (all reals), the \emph{codomain} is the declared target (all reals), and the \emph{range} is the actual list of outputs ($[0,\infty)$). This function is neither injective (since $f(2)=f(-2)$) nor surjective onto $\R$ (negative numbers cannot occur). If the domain is restricted to $[0,\infty)$ and the codomain to $[0,\infty)$, the same formula becomes bijective. Thus, \emph{do not discuss injectivity or surjectivity without looking at the exact domain and codomain}.
\subsection{Direct image versus inverse image}
For $f(x)=x^2$ on $\R$, the direct image of $[-2,-1]$ is $[1,4]$. The inverse image of $[1,4]$ is $[-2,-1]\cup[1,2]$: you ask for \emph{all inputs whose outputs land in the target set}. The inverse image exists even though the function itself has no inverse. Images need not preserve intersections because different inputs can collapse to the same output; inverse images always preserve intersections because membership is checked on the same input.
\subsection{Equivalence relations as grouping rules}
An equivalence relation is a defensible definition of ``belongs in the same group'': every object belongs with itself (reflexive); if $a$ groups with $b$ then $b$ groups with $a$ (symmetric); and grouping must chain from $a$ via $b$ to $c$ (transitive). Congruence modulo $3$ groups integers by remainder: $[0]=\{\ldots,-3,0,3,\ldots\}$, $[1]=\{\ldots,-2,1,4,\ldots\}$, and $[2]=\{\ldots,-1,2,5,\ldots\}$. These classes partition $\Z$ into disjoint nonempty blocks.
\subsection{Why countability is not the same as finiteness}
A set is countable when its elements can be arranged into a finite or natural-number-indexed list. Both $\N$ and $\Z$ are infinite, but $\Z$ is listable as $0,1,-1,2,-2,\ldots$. Rational numbers are also countable, despite their density on the real line. Real numbers cannot be listed this way. Cantor's theorem formalises the key obstacle: the power set $\Pow(S)$ always has strictly larger cardinality than $S$.
\begin{keybox}{A miniature proof of Cantor's theorem}
Suppose a proposed surjection $f:S\to\Pow(S)$ lists all subsets of $S$. Form the diagonal subset $D=\{s\in S:s\notin f(s)\}$. If $f(a)=D$, then $a\in D\iff a\notin f(a)=D$, an impossibility. No surjection can exist.\end{keybox}
\subsection{Partial orders: a frequent conceptual trap}
A partial order need not compare every pair. For example, inclusion $\subseteq$ is a partial order on $\Pow(\{1,2\})$, but $\{1\}$ and $\{2\}$ are incomparable. In a \emph{total} order every pair is comparable. A \emph{maximal} element has no strictly larger comparable element; a \emph{maximum} is at least as large as every element. These words are not interchangeable.
\newpage\section{Examination writing: model structures and common pitfalls}
\begin{description}[style=nextline,leftmargin=13mm]
\item[Prove a distance is a metric] State the set and the proposed map, then test MS0, MS1a, MS1b, MS2, MS3 in order. When the question asks \emph{only} for MS3, do not spend most of your answer reproving the other axioms.
\item[Prove a set identity] Write $x\in$ left-hand side, translate symbols into logic (``and'', ``or'', ``not'', ``exists'', ``for every''), then reconstruct the right-hand side.
\item[Prove injectivity] Assume $f(a)=f(b)$ and deduce $a=b$. To disprove it, give distinct $a,b$ with equal values.
\item[Prove surjectivity] Take arbitrary $y$ in the declared codomain and solve $f(x)=y$ with $x$ in the domain. To disprove it, exhibit a target value not achieved.
\item[Show a map is bijective] An explicit inverse with both compositions verified is particularly clear, but an independent injectivity and surjectivity argument also suffices.
\item[Count a collection] Explain why the enumeration or encoding is exhaustive and why duplicates are impossible when claiming a bijection.
\item[Handle an infimum] Work with admissible numbers within an arbitrary $\varepsilon>0$ of the infimum; do not automatically assume the infimum is attained.
\end{description}
\begin{tipbox}{Marks: important distinction}The source handouts print a mark total for some exercises, but not for all Section 1 questions. Every breakdown below is an \emph{illustrative, reasonable grading allocation}, \textbf{not an official university marking scheme}. Mark allocations for source-labelled questions add up to the printed total. For questions with no printed total, a clearly labelled suggested 10-mark scale is used.\end{tipbox}

\newpage\section{Complete worked solutions: Section 0}
\ex{0.1}{Discrete metric (10 marks)}
Define $d(s,t)=0$ when $s=t$, and $1$ otherwise. \textbf{MS0:} values lie in $\{0,1\}\subset[0,\infty)$, so are finite and nonnegative. \textbf{MS1:} by its definition $d(s,t)=0$ precisely when $s=t$. \textbf{MS2:} $s=t$ iff $t=s$, so the same case applies to $d(s,t)$ and $d(t,s)$. \textbf{MS3:} let $s,t,u\in S$. If $s=u$, the left side is $0$ and the inequality follows. If $s\ne u$, $d(s,u)=1$. At least one of $s\ne t$ or $t\ne u$ must hold: otherwise $s=t=u$, contradiction. Thus at least one of $d(s,t),d(t,u)$ equals $1$, giving $1\le d(s,t)+d(t,u)$. All axioms hold. $\square$

\begin{tipbox}{Source total: illustrative mark-by-mark checklist}
\begin{tabular}{@{}p{0.77\linewidth}r@{}}
Correct codomain and non-negativity & 1 marks \\
Identity of indiscernibles & 2 marks \\
Symmetry & 2 marks \\
Triangle inequality with both cases & 4 marks \\
Explicit conclusion & 1 marks \\
\hline\textbf{Total} & \textbf{10 marks}\\
\end{tabular}
\end{tipbox}
\ex{0.2}{Integration metric: definiteness (10 marks)}
Let $h\in C([0,1];\R)$, and suppose $h\not\equiv0$. Choose $t_0\in[0,1]$ with $|h(t_0)|=a>0$. By continuity at $t_0$, there exists $\varepsilon>0$ such that $|h(t)-h(t_0)|<a/2$ whenever $t\in[0,1]$ and $|t-t_0|<\varepsilon$. The reverse triangle inequality yields $|h(t)|\ge |h(t_0)|-|h(t)-h(t_0)|>a/2$ throughout this neighbourhood. Its intersection with $[0,1]$ contains an interval of positive length $\ell>0$ (including if $t_0$ is an endpoint). Consequently
\[\int_0^1|h(t)|\,dt\ge\int_{\text{that interval}}|h(t)|\,dt\ge\frac a2\ell>0.\]
This proves the contrapositive: if $\int_0^1|h|=0$, then $h\equiv0$. Taking $h=f-g$ proves $d_1(f,g)=0\Rightarrow f=g$. \textbf{Why continuity matters:} a function equal to $1$ at one point and $0$ elsewhere would have zero integral without being identically zero. $\square$

\begin{tipbox}{Source total: illustrative mark-by-mark checklist}
\begin{tabular}{@{}p{0.77\linewidth}r@{}}
Choose a nonzero point & 2 marks \\
Apply continuity to get positive lower bound & 3 marks \\
Integrate over interval of positive length & 3 marks \\
Contrapositive and definitive conclusion & 2 marks \\
\hline\textbf{Total} & \textbf{10 marks}\\
\end{tabular}
\end{tipbox}
\ex{0.3}{Supremum metric: triangle and definiteness (10 marks)}
Write $d_\infty(f,g)=\sup_{t\in I}|f(t)-g(t)|$ with $I=[0,1]$. \textbf{MS1b:} if $d_\infty(f,g)=0$, for every $t$ we have $0\le|f(t)-g(t)|\le0$, hence $f(t)=g(t)$ for every $t$, i.e. $f=g$. \textbf{MS3:} for each $t\in I$,
\[|f(t)-h(t)|\le |f(t)-g(t)|+|g(t)-h(t)|\le d_\infty(f,g)+d_\infty(g,h).\]
Taking the supremum on the left (right side is independent of $t$) yields $d_\infty(f,h)\le d_\infty(f,g)+d_\infty(g,h)$. $\square$

\begin{tipbox}{Source total: illustrative mark-by-mark checklist}
\begin{tabular}{@{}p{0.77\linewidth}r@{}}
Derive MS1b pointwise & 4 marks \\
Pointwise triangle inequality and supremum & 4 marks \\
Conclude both axioms & 2 marks \\
\hline\textbf{Total} & \textbf{10 marks}\\
\end{tabular}
\end{tipbox}
\ex{0.4}{Counting metric: triangle inequality (10 marks)}
For $A,B,C\subseteq S$, the crucial set inclusion is
\[A\triangle C\subseteq(A\triangle B)\cup(B\triangle C).\]
To justify it, take $x\in A\triangle C$: its membership in $A$ and $C$ differs. If membership in $A$ also differs from that in $B$, then $x\in A\triangle B$. Otherwise its membership in $A$ and $B$ agrees, so membership in $B$ and $C$ differs, giving $x\in B\triangle C$. Because $S$ is finite, cardinalities obey
\begin{align*}d(A,C)&=|A\triangle C|\\&\le|(A\triangle B)\cup(B\triangle C)|\\&\le|A\triangle B|+|B\triangle C|\\&=d(A,B)+d(B,C).\end{align*}
Hence the counting distance satisfies MS3. \textbf{Intuition:} each change from $A$ to $C$ must occur either on the $A$--$B$ leg or the $B$--$C$ leg. $\square$

\begin{tipbox}{Source total: illustrative mark-by-mark checklist}
\begin{tabular}{@{}p{0.77\linewidth}r@{}}
Prove symmetric-difference inclusion & 4 marks \\
Bound cardinalities and identify distances & 4 marks \\
State triangle conclusion & 2 marks \\
\hline\textbf{Total} & \textbf{10 marks}\\
\end{tabular}
\end{tipbox}
\ex{0.5}{Two unusual metrics on $\R^2$ (20 marks)}
Set $r(u)=\|u\|_2$. For the first metric, if $u=w$, $d(u,w)=0\le d(u,v)+d(v,w)$. If $u\ne w$, there are three cases. If $v=u$, then $d(u,v)=0$ and $d(v,w)=d(u,w)$. If $v=w$, similarly equality holds. Otherwise both $u\ne v$ and $v\ne w$, and
\[d(u,v)+d(v,w)=(r(u)+r(v))+(r(v)+r(w))\ge r(u)+r(w)=d(u,w).\]
Thus $d_{\rm TGV}$ satisfies MS3.

For $d'_{\rm TGV}$, observe first that
\[d'(a,b)\ge\|a-b\|_2\quad\text{for all }a,b,\]
because on dependent pairs it is equality, and on independent pairs this is the norm triangle inequality. If $u,w$ are dependent, then
\[d'(u,w)=\|u-w\|_2\le\|u-v\|_2+\|v-w\|_2\le d'(u,v)+d'(v,w).\]
If $u,w$ are independent, then $u,w\ne0$ and $d'(u,w)=r(u)+r(w)$. If $v=0$, the two legs have lengths $r(u)$ and $r(w)$, so equality holds. If $v\ne0$, it cannot be dependent on \emph{both} independent vectors $u,w$. If $u,v$ are independent, $d'(u,v)=r(u)+r(v)$, while $d'(v,w)\ge\|v-w\|\ge r(w)-r(v)$, so their sum is at least $r(u)+r(w)$. Otherwise $u,v$ are dependent, hence $v,w$ independent, and $d'(v,w)=r(v)+r(w)$ while $d'(u,v)\ge r(u)-r(v)$; again the total is at least $r(u)+r(w)$. All cases prove MS3. $\square$

\begin{tipbox}{Source total: illustrative mark-by-mark checklist}
\begin{tabular}{@{}p{0.77\linewidth}r@{}}
First metric: zero/end-point cases & 3 marks \\
First metric: remaining case and conclusion & 5 marks \\
Second metric: Euclidean lower bound & 3 marks \\
Dependent endpoints case & 3 marks \\
Independent endpoints case including subcases & 5 marks \\
Conclusion for both metrics & 1 marks \\
\hline\textbf{Total} & \textbf{20 marks}\\
\end{tabular}
\end{tipbox}
\ex{0.6}{Hausdorff metric: triangle inequality (10 marks)}
\textbf{(a)(i)} If $A\subseteq B$ and $x\in A^\delta$, choose $a\in A$ with $\|x-a\|\le\delta$. Since $a\in B$, $x\in B^\delta$. Therefore $A^\delta\subseteq B^\delta$.

\textbf{(a)(ii)} If $x\in(A^\gamma)^\delta$, choose $y\in A^\gamma$ with $\|x-y\|\le\delta$. Choose $a\in A$ with $\|y-a\|\le\gamma$. Euclidean MS3 gives $\|x-a\|\le\delta+\gamma$, hence $x\in A^{\gamma+\delta}$. Thus $(A^\gamma)^\delta\subseteq A^{\gamma+\delta}$.

\textbf{(b)} Let $K,J,L$ be nonempty bounded closed sets and write $\alpha=d_H(K,J)$, $\beta=d_H(J,L)$. Given any $\varepsilon>0$, the definition of infimum permits admissible radii $\gamma<\alpha+\varepsilon/2$ and $\delta<\beta+\varepsilon/2$ so that $K^\gamma\supseteq J$, $J^\gamma\supseteq K$, $J^\delta\supseteq L$ and $L^\delta\supseteq J$. Applying (a),
\[L\subseteq J^\delta\subseteq(K^\gamma)^\delta\subseteq K^{\gamma+\delta},\]
and likewise $K\subseteq L^{\gamma+\delta}$. Thus $d_H(K,L)\le\gamma+\delta<\alpha+\beta+\varepsilon$. As $\varepsilon>0$ was arbitrary, $d_H(K,L)\le d_H(K,J)+d_H(J,L)$. Together with the other properties established in the notes, this proves the metric axioms. $\square$
\newpage\section{Complete worked solutions: Section 1}

\begin{tipbox}{Source total: illustrative mark-by-mark checklist}
\begin{tabular}{@{}p{0.77\linewidth}r@{}}
Monotonicity of fattening & 2 marks \\
Successive-fattening inclusion & 3 marks \\
Admissible radii and two inclusions & 3 marks \\
Epsilon limit and conclusion & 2 marks \\
\hline\textbf{Total} & \textbf{10 marks}\\
\end{tabular}
\end{tipbox}
\ex{1.1}{Subset equivalences and symmetric difference}
\textbf{(a)} Assume $A\subseteq B$. Then every member of $A$ belongs to both $A$ and $B$, so $A\subseteq A\cap B$; conversely $A\cap B\subseteq A$, hence $A\cap B=A$. Conversely, if $A\cap B=A$, then each $x\in A$ also lies in $B$, giving $A\subseteq B$. Similarly $A\subseteq B\Rightarrow A\cup B=B$ because adjoining members of $A$ adds nothing to $B$. If $A\cup B=B$, every $a\in A$ lies in $B$, hence $A\subseteq B$. Thus all three assertions are equivalent.

\textbf{(b)} For any $x$, membership in $(A\cup B)\setminus(A\cap B)$ means $x$ belongs to at least one set but not both, exactly equivalent to $(x\in A\land x\notin B)$ or $(x\in B\land x\notin A)$. This is membership in $(A\setminus B)\cup(B\setminus A)$. \textbf{(c)} Draw two overlapping circles: shade only the two outer crescent regions, leaving the overlap unshaded. $\square$

\begin{tipbox}{Suggested practice total: mark-by-mark checklist}
\begin{tabular}{@{}p{0.77\linewidth}r@{}}
Establish equivalence of three statements & 4 marks \\
Prove symmetric-difference identity & 4 marks \\
Correct shaded-region description & 2 marks \\
\hline\textbf{Total} & \textbf{10 marks}\\
\end{tabular}
\end{tipbox}
\ex{1.2}{Distributivity and De Morgan}
Let $x\in S$. \textbf{(a)}
\begin{align*}x\in A\cap\bigcup\mathcal C&\iff x\in A\text{ and }(\exists C\in\mathcal C)\,x\in C\\&\iff(\exists C\in\mathcal C)\,x\in A\cap C\\&\iff x\in\bigcup_{C\in\mathcal C}(A\cap C).\end{align*}
\textbf{(b)}
\begin{align*}x\in A\setminus\bigcap\mathcal C&\iff x\in A\text{ and not }[(\forall C\in\mathcal C)\,x\in C]\\&\iff(\exists C\in\mathcal C)\,[x\in A\text{ and }x\notin C]\\&\iff x\in\bigcup_{C\in\mathcal C}(A\setminus C).\end{align*}
Both equivalences hold for an empty family with the stated universe conventions. $\square$

\begin{tipbox}{Suggested practice total: mark-by-mark checklist}
\begin{tabular}{@{}p{0.77\linewidth}r@{}}
Distributive membership proof & 5 marks \\
De Morgan membership proof & 5 marks \\
\hline\textbf{Total} & \textbf{10 marks}\\
\end{tabular}
\end{tipbox}
\ex{1.3}{Counting injections and surjections}
For $S=\{1,2\}$, $T=\{1,2,3\}$: \textbf{(a)} choose $f(1)$ in $3$ ways; to avoid collision choose $f(2)$ among the remaining $2$. Answer $3\cdot2=\boxed6$. \textbf{(b)} A function from a two-element domain has at most two different outputs and cannot cover three targets: $\boxed0$. \textbf{(c)} There are $2^3=8$ total maps $T\to S$; exactly two are not onto (both constant maps). Hence $\boxed6$. Alternatively choose the nonempty preimage of $1$: any of $2^3-2=6$ proper, nonempty subsets. $\square$

\begin{tipbox}{Suggested practice total: mark-by-mark checklist}
\begin{tabular}{@{}p{0.77\linewidth}r@{}}
Count injections & 4 marks \\
Explain impossibility of surjections S to T & 2 marks \\
Count surjections T to S & 4 marks \\
\hline\textbf{Total} & \textbf{10 marks}\\
\end{tabular}
\end{tipbox}
\ex{1.4}{Composition preserves injectivity and surjectivity}
\textbf{(a)(i)} If $(g\circ f)(x)=(g\circ f)(y)$, then $g(f(x))=g(f(y))$. Injectivity of $g$ gives $f(x)=f(y)$; injectivity of $f$ then gives $x=y$. \textbf{(a)(ii)} If $f(x)=f(y)$ then $(g\circ f)(x)=(g\circ f)(y)$; injectivity of the composition yields $x=y$.

\textbf{(b)(i)} For any $u\in U$, surjectivity of $g$ supplies $t\in T$ with $g(t)=u$. Surjectivity of $f$ supplies $s\in S$ with $f(s)=t$. Thus $(g\circ f)(s)=u$. \textbf{(b)(ii)} If $g\circ f$ is onto and $u\in U$, find $s$ with $g(f(s))=u$ and set $t=f(s)$, proving $g$ is onto.

\textbf{(c)(i)} Parts (a)(i), (b)(i) say the composition is both injective and surjective. \textbf{(c)(ii)} Parts (a)(ii), (b)(ii) give the necessary conditions. \textbf{Warning:} $g\circ f$ bijective need not make both $f$ and $g$ bijective. $\square$

\begin{tipbox}{Suggested practice total: mark-by-mark checklist}
\begin{tabular}{@{}p{0.77\linewidth}r@{}}
Composition preserves injectivity & 2 marks \\
Infer injectivity of first map & 2 marks \\
Composition preserves surjectivity & 2 marks \\
Infer surjectivity of second map & 2 marks \\
Bijective implications & 2 marks \\
\hline\textbf{Total} & \textbf{10 marks}\\
\end{tabular}
\end{tipbox}
\ex{1.5}{Injectivity of a rational function}
Let $f:\R\setminus\{1\}\to\R$ be $f(x)=2x/(x-1)$. \textbf{(a)} If $2x/(x-1)=2y/(y-1)$, cross multiplication (denominators nonzero) gives $x(y-1)=y(x-1)$, hence $xy-x=xy-y$ and $x=y$. Thus $f$ injective.

\textbf{(b)} Put $t=2x/(x-1)$. Rearrangement gives $t(x-1)=2x$, so $(t-2)x=t$. If $t=2$ this would read $0=2$, impossible. For any $t\ne2$, $x=t/(t-2)$ is well-defined and differs from $1$ (otherwise $t=t-2$). So $\im f=\boxed{\R\setminus\{2\}}$.

\textbf{(c)} On $t\ne2$ the inverse formula is $t/(t-2)$. To make a left inverse with domain \emph{all} of $\R$, assign any permissible value at $2$, say $0$:
\[g(t)=\begin{cases}\dfrac{t}{t-2},&t\ne2,\\0,&t=2.\end{cases}\]
Then $g:\R\to\R\setminus\{1\}$, and $g(f(x))=(2x/(x-1))/(2x/(x-1)-2)=x$. Thus $g\circ f=\id_{\R\setminus\{1\}}$. $\square$

\begin{tipbox}{Source total: illustrative mark-by-mark checklist}
\begin{tabular}{@{}p{0.77\linewidth}r@{}}
Injectivity algebra & 3 marks \\
Determine the range & 3 marks \\
Define left inverse on entire codomain & 3 marks \\
Verify composition & 1 marks \\
\hline\textbf{Total} & \textbf{10 marks}\\
\end{tabular}
\end{tipbox}
\ex{1.6}{Functions as indexed products}
Let $P=\prod_{s\in S}T_s$ with $T_s=T$, and $F=\mathcal F(S;T)$. Define $\phi:P\to F$ by $\phi((t_s)_{s\in S})(s)=t_s$. Define the proposed inverse $\psi:F\to P$ by $\psi(f)=(f(s))_{s\in S}$. This is well-defined since each $f(s)\in T=T_s$. For $p=(t_s)\in P$, $\psi(\phi(p))=(\phi(p)(s))_{s\in S}=(t_s)_{s\in S}=p$. For $f\in F$, $\phi(\psi(f))(s)=f(s)$ for all $s$, hence $\phi(\psi(f))=f$. Both compositions are identities, so $\phi$ is bijective. If $S=\varnothing$, both sets have one element; if $S\ne\varnothing,T=\varnothing$, both sets are empty. The argument remains valid in each case. $\square$

\begin{tipbox}{Suggested practice total: mark-by-mark checklist}
\begin{tabular}{@{}p{0.77\linewidth}r@{}}
Construct inverse on functions & 4 marks \\
Verify both compositions & 4 marks \\
Handle special cases & 2 marks \\
\hline\textbf{Total} & \textbf{10 marks}\\
\end{tabular}
\end{tipbox}
\ex{1.7}{Power sets and indicator functions}
Define $\chi:\Pow(S)\to\mathcal F(S;\{0,1\})$ by $\chi(A)=1_A$. Given $h:S\to\{0,1\}$ define $\sigma(h)=\{s\in S:h(s)=1\}\in\Pow(S)$. Then $\sigma(\chi(A))=\{s:1_A(s)=1\}=A$. Conversely, for every $h$ and $s$, $1_{\sigma(h)}(s)=1$ iff $h(s)=1$, and is $0$ iff $h(s)=0$. Because $h(s)$ has only these two possible values, $\chi(\sigma(h))=h$. Thus $\sigma=\chi^{-1}$ and the sets are bijective. $\square$

\begin{tipbox}{Source total: illustrative mark-by-mark checklist}
\begin{tabular}{@{}p{0.77\linewidth}r@{}}
Define inverse via the set of ones & 4 marks \\
Verify sigma after chi & 3 marks \\
Verify chi after sigma & 3 marks \\
\hline\textbf{Total} & \textbf{10 marks}\\
\end{tabular}
\end{tipbox}
\ex{1.8}{Images of unions and intersections}
Let $\mathcal A$ be a family of subsets of $S$. \textbf{(a)} For $t\in T$,
\begin{align*}t\in f[\bigcup\mathcal A]&\iff(\exists s\in\bigcup\mathcal A)\ f(s)=t\\&\iff(\exists A\in\mathcal A)(\exists s\in A)\ f(s)=t\\&\iff t\in\bigcup_{A\in\mathcal A}f[A].\end{align*}
\textbf{(b)} If $t\in f[\bigcap\mathcal A]$, choose $s\in\bigcap\mathcal A$ with $f(s)=t$. The same $s$ is in every $A$, so $t\in f[A]$ for every $A$, hence $t\in\bigcap_A f[A]$. It need not be equality when $f$ is not injective.

\textbf{Additional identity from the lecture slides:} If $t\in f[A]\setminus f[B]$, choose $a\in A$ with $f(a)=t$. If $a\in B$ then $t\in f[B]$, contradiction. Thus $a\in A\setminus B$ and $t\in f[A\setminus B]$, proving $f[A\setminus B]\supseteq f[A]\setminus f[B]$. Equality need not hold, e.g. a constant map on distinct $a,b$, with $A=\{a\},B=\{b\}$. $\square$

\begin{tipbox}{Suggested practice total: mark-by-mark checklist}
\begin{tabular}{@{}p{0.77\linewidth}r@{}}
Union image identity & 5 marks \\
Intersection image inclusion & 5 marks \\
\hline\textbf{Total} & \textbf{10 marks}\\
\end{tabular}
\end{tipbox}
\ex{1.9}{Inverse images of unions and differences}
For any $s\in S$,
\begin{align*}s\in f^{-1}[\bigcup\mathcal C]&\iff f(s)\in\bigcup\mathcal C\\&\iff(\exists C\in\mathcal C)\ f(s)\in C\\&\iff s\in\bigcup_{C\in\mathcal C}f^{-1}[C].\end{align*}
Likewise
\begin{align*}s\in f^{-1}[D\setminus C]&\iff f(s)\in D\text{ and }f(s)\notin C\\&\iff s\in f^{-1}[D]\text{ and }s\notin f^{-1}[C]\\&\iff s\in f^{-1}[D]\setminus f^{-1}[C].\end{align*}
Since $f^{-1}[T]=S$, taking $D=T$ gives $f^{-1}[T\setminus C]=S\setminus f^{-1}[C]$. \textbf{Further lecture-slide identity:} $f^{-1}[\bigcap\mathcal C]=\bigcap_{C\in\mathcal C}f^{-1}[C]$, by replacing ``there exists'' by ``for every'' in the first proof. $\square$

\begin{tipbox}{Source total: illustrative mark-by-mark checklist}
\begin{tabular}{@{}p{0.77\linewidth}r@{}}
Preimage of union using membership & 4 marks \\
Preimage of difference using membership & 4 marks \\
Complement specialisation & 2 marks \\
\hline\textbf{Total} & \textbf{10 marks}\\
\end{tabular}
\end{tipbox}
\ex{1.10}{Disjoint union of equivalence relations}
Let $S=A_1\cup A_2$ with $A_1\cap A_2=\varnothing$. Define $x\sim y$ if they belong to the same $A_i$ and $x\sim_i y$. \textbf{Reflexivity:} any $x$ lies in exactly one $A_i$, so $x\sim_i x$ and therefore $x\sim x$. \textbf{Symmetry:} if $x\sim y$, then $x,y\in A_i$ and $x\sim_i y$ for some $i$. Symmetry of $\sim_i$ yields $y\sim_i x$, hence $y\sim x$. \textbf{Transitivity:} if $x\sim y$ and $y\sim z$, then the first pair is in some $A_i$ and the second in some $A_j$. Since the shared $y$ belongs to both and the sets are disjoint, $i=j$. Transitivity in $A_i$ gives $x\sim_i z$, hence $x\sim z$. All three properties hold. $\square$

\begin{tipbox}{Suggested practice total: mark-by-mark checklist}
\begin{tabular}{@{}p{0.77\linewidth}r@{}}
Reflexivity & 3 marks \\
Symmetry & 3 marks \\
Transitivity using disjointness & 4 marks \\
\hline\textbf{Total} & \textbf{10 marks}\\
\end{tabular}
\end{tipbox}
\ex{1.11}{Symmetry and transitivity do not imply reflexivity}
\textbf{(a)} The argument silently assumes the existence of a $y$ with $x\sim y$ for each $x$. Symmetry can reverse a relation \emph{only if one already holds}; it cannot produce a pair relating a completely unrelated element. Therefore the step ``take $y$ such that $x\sim y$'' may have no witness.

\textbf{(b)} Take nonempty $S=\{1\}$ with $\sim=\varnothing$ (the empty binary relation). Symmetry and transitivity are vacuously true because there are no pairs $(x,y)$ in the relation and no applicable chains. Reflexivity fails, since $1\not\sim1$. $\square$

\begin{tipbox}{Source total: illustrative mark-by-mark checklist}
\begin{tabular}{@{}p{0.77\linewidth}r@{}}
Identify missing existential assumption & 5 marks \\
Valid symmetric and transitive counterexample & 5 marks \\
\hline\textbf{Total} & \textbf{10 marks}\\
\end{tabular}
\end{tipbox}
\ex{1.12}{Explicit bijection from $\N$ to $\Z$}
Assume $\N=\{1,2,3,\ldots\}$. Enumerate $\Z$ as $0,1,-1,2,-2,3,-3,\ldots$ via
\[f(1)=0,\qquad f(2k)=k,\qquad f(2k+1)=-k\quad(k\ge1).\]
Every $z\in\Z$ is hit: $z=0$ by input $1$, positive $z=k$ by $2k$, and negative $z=-k$ by $2k+1$. Different positive even inputs give distinct positive outputs, distinct odd inputs beyond $1$ give distinct negative outputs, and $1$ uniquely gives $0$. So no outputs collide. Hence $f$ is a bijection, proving $\Z$ countably infinite. $\square$

\begin{tipbox}{Suggested practice total: mark-by-mark checklist}
\begin{tabular}{@{}p{0.77\linewidth}r@{}}
Give explicit enumeration/function & 4 marks \\
Prove surjectivity & 3 marks \\
Prove injectivity & 3 marks \\
\hline\textbf{Total} & \textbf{10 marks}\\
\end{tabular}
\end{tipbox}
\ex{1.13}{Rational-coefficient polynomials are countable}
For each $n\ge0$, a polynomial of degree at most $n$ is uniquely represented by its coefficient vector $(a_0,\ldots,a_n)\in\Q^{n+1}$ through $\sum_{k=0}^n a_kx^k$. Thus the set $S_n$ of such polynomials is in bijection with $\Q^{n+1}$. Since $\Q$ is countable and finite Cartesian products of countable sets are countable, each $S_n$ is countable. Every polynomial has some finite degree (the zero polynomial also lies in every $S_n$), so $S=\bigcup_{n=0}^{\infty}S_n$. A countable union of countable sets is countable; hence $S$ is countable. \textbf{Pitfall:} ``infinitely many possible coefficients'' does not mean uncountable: each individual polynomial uses only finitely many coefficients. $\square$

\begin{tipbox}{Suggested practice total: mark-by-mark checklist}
\begin{tabular}{@{}p{0.77\linewidth}r@{}}
Express bounded-degree polynomials via rational tuples & 4 marks \\
Use finite-product countability & 3 marks \\
Use countable union & 3 marks \\
\hline\textbf{Total} & \textbf{10 marks}\\
\end{tabular}
\end{tipbox}
\ex{1.14}{Finite subsets of $\N$ are countable}
\textbf{First proof (binary encoding, especially elegant).} Define
\[E:\Pow_{00}(\N)\to\N_0,\qquad E(A)=\sum_{n\in A}2^{n-1},\]
where $\N_0=\{0,1,2,\ldots\}$ and the empty sum equals zero. Every $A$ is finite so the sum is defined. Unique binary expansion implies $E(A)=E(B)$ only if $A=B$. Thus $E$ is injective and $\Pow_{00}(\N)$ is countable. In fact every nonnegative integer has a unique finite binary expansion, so $E$ is bijective.

\textbf{Proof following the lecture hint.} Decompose $\Pow_{00}(\N)=\bigcup_{k\ge0}\Pow^{(k)}(\N)$ by cardinality. Associate to each $k$-element subset its uniquely increasing $k$-tuple $(a_1<\cdots<a_k)$ in $\N^k$. This injects $\Pow^{(k)}(\N)$ into the countable set $\N^k$, so each layer is countable. Their countable union is countable. $\square$

\begin{tipbox}{Suggested practice total: mark-by-mark checklist}
\begin{tabular}{@{}p{0.77\linewidth}r@{}}
Define finite-subset encoding & 4 marks \\
Prove uniqueness or injection & 4 marks \\
Conclude countability & 2 marks \\
\hline\textbf{Total} & \textbf{10 marks}\\
\end{tabular}
\end{tipbox}
\ex{1.15}{Total versus strict total orders}
\textbf{(a)} From a total order $\le$, define $s<t$ iff $s\le t$ and $s\ne t$. \emph{Trichotomy:} totality gives $s\le t$ or $t\le s$; if $s\ne t$, antisymmetry prevents both, so exactly one strict direction holds. If equal, neither strict direction holds. \emph{Transitivity:} $s<t$ and $t<u$ give $s\le t\le u$, hence $s\le u$; if $s=u$, antisymmetry applied along the chain forces $s=t$, contradiction. Therefore $s<u$.

\textbf{(b)} Conversely define $s\le t$ iff $s<t$ or $s=t$. \emph{Reflexivity:} $s=s$, hence $s\le s$. \emph{Antisymmetry:} if $s\le t$ and $t\le s$ with $s\ne t$, then $s<t$ and $t<s$, contradicting trichotomy. \emph{Transitivity:} if $s\le t\le u$, cases involving equality are immediate; otherwise $s<t$ and $t<u$, hence $s<u$ by strict transitivity. \emph{Totality:} trichotomy gives one of $s<t$, $s=t$, $t<s$, ensuring $s\le t$ or $t\le s$. Thus $\le$ is a total order. $\square$

\begin{tipbox}{Source total: illustrative mark-by-mark checklist}
\begin{tabular}{@{}p{0.77\linewidth}r@{}}
Trichotomy from totality and antisymmetry & 5 marks \\
Transitivity of strict relation & 5 marks \\
Reflexivity of non-strict relation & 2 marks \\
Antisymmetry & 3 marks \\
Transitivity & 3 marks \\
Totality & 2 marks \\
\hline\textbf{Total} & \textbf{20 marks}\\
\end{tabular}
\end{tipbox}
\ex{1.16}{Bijections of two three-element sets}
A bijection is specified by the ordered triple $(f(x),f(y),f(z))$, a permutation of $(a,b,c)$. Evaluate statements $A:f(x)=a$, $B:f(y)\ne a$, and $C:f(z)\ne b$:
\begin{center}\begin{tabular}{c|ccc|c}$(f(x),f(y),f(z))$&A&B&C&true count\\\hline $(a,b,c)$&T&T&T&3\\$(a,c,b)$&T&T&F&2\\$(b,a,c)$&F&F&T&1\\$(b,c,a)$&F&T&T&2\\$(c,a,b)$&F&F&F&0\\$(c,b,a)$&F&T&T&2\end{tabular}\end{center}
\textbf{(a)} Exactly two true: $\boxed{(a,c,b),\ (b,c,a),\ (c,b,a)}$. \textbf{(b)} Exactly one true: $\boxed{(b,a,c)}$. This exhaustive table works because there are $3!=6$ bijections. $\square$

\begin{tipbox}{Suggested practice total: mark-by-mark checklist}
\begin{tabular}{@{}p{0.77\linewidth}r@{}}
List all six bijections & 4 marks \\
Evaluate truth conditions correctly & 4 marks \\
Identify requested cases & 2 marks \\
\hline\textbf{Total} & \textbf{10 marks}\\
\end{tabular}
\end{tipbox}
\ex{1.17}{A relation defined from a function and partition}
The notes define $t_1\sim t_2$ iff there is a block $A\in\mathcal C$ such that $f^{-1}[\{t_1,t_2\}]\subseteq A$, where $\mathcal C$ partitions $S$.

\textbf{(a) Reflexivity: not necessarily.} Let $S=\{p,q\}$, $T=\{t\}$, $f(p)=f(q)=t$, and partition $\mathcal C=\{\{p\},\{q\}\}$. The fibre $f^{-1}[\{t\}]=\{p,q\}$ lies in no single block, so $t\not\sim t$. \emph{If $f$ is injective, yes:} each fibre $f^{-1}[\{t\}]$ is empty or a singleton, and any singleton lies in some block. Provided the partition of nonempty $S$ has a block, the empty fibre is also a subset of a block. Hence $t\sim t$ for all $t$.

\textbf{(b) Symmetry: always.} Since $\{t_1,t_2\}=\{t_2,t_1\}$, their inverse images coincide. Thus $t_1\sim t_2$ iff $t_2\sim t_1$, with or without bijectivity.

\textbf{(c) Transitivity: not necessarily.} Let $S=\{p,q\}$, $T=\{a,b,c\}$, $f(p)=a, f(q)=c$, and $\mathcal C=\{\{p\},\{q\}\}$. Since $b\notin\im f$, its fibre is empty. We have $f^{-1}[\{a,b\}]=\{p\}$ and $f^{-1}[\{b,c\}]=\{q\}$, so $a\sim b$ and $b\sim c$. But $f^{-1}[\{a,c\}]=\{p,q\}$ is not contained in any block, so $a\not\sim c$.

\emph{If $f$ is surjective, transitivity does hold.} Suppose $t_1\sim t_2$ and $t_2\sim t_3$. Choose blocks $A,B$ with $f^{-1}[\{t_1,t_2\}]\subseteq A$ and $f^{-1}[\{t_2,t_3\}]\subseteq B$. Surjectivity makes $f^{-1}[\{t_2\}]$ nonempty, so a point in that fibre belongs to both blocks. Since blocks in a partition are disjoint unless equal, $A=B$. Therefore $f^{-1}[\{t_1,t_3\}]\subseteq A$ and $t_1\sim t_3$. $\square$

\begin{tipbox}{Suggested practice total: mark-by-mark checklist}
\begin{tabular}{@{}p{0.77\linewidth}r@{}}
Reflexivity analysis and conditions & 3 marks \\
Symmetry & 2 marks \\
Transitivity counterexample and sufficient condition & 5 marks \\
\hline\textbf{Total} & \textbf{10 marks}\\
\end{tabular}
\end{tipbox}
\ex{1.18}{Real numbers are uncountable}
\textbf{(a)} Define $\psi:\{0,1\}^{\N}\to\R$ by $\psi(a)=\sum_{n=1}^{\infty}a_n10^{-n}$. It converges absolutely since $0\le\psi(a)\le\sum 10^{-n}=1/9$. Let $a\ne b$ and let $k$ be their first differing index. Without loss, $a_k=1$, $b_k=0$. Their difference is
\[\psi(a)-\psi(b)=10^{-k}+\sum_{n>k}(a_n-b_n)10^{-n}\ge10^{-k}-\sum_{n>k}10^{-n}.\]
The tail is $\sum_{n>k}10^{-n}=10^{-k}/9$, so $\psi(a)-\psi(b)\ge(8/9)10^{-k}>0$. Thus $\psi$ is injective. The use of digits \emph{0 and 1 only} avoids ambiguous decimal expansions.

\textbf{(b)} The set $\{0,1\}^{\N}$ is bijective to $\mathcal F(\N;\{0,1\})$, which is bijective to $\Pow(\N)$ via indicator functions. By Cantor's theorem $\Pow(\N)$ is uncountable, hence so is $\{0,1\}^{\N}$. If $\R$ were countable, every subset of it, including the image of the injection $\psi$, would be countable, contradiction. So $\boxed{\R\text{ is uncountable}}$. $\square$

\begin{tipbox}{Suggested practice total: mark-by-mark checklist}
\begin{tabular}{@{}p{0.77\linewidth}r@{}}
Define convergent digit map & 4 marks \\
Prove injectivity via first differing digit & 3 marks \\
Invoke Cantor to establish uncountability & 3 marks \\
\hline\textbf{Total} & \textbf{10 marks}\\
\end{tabular}
\end{tipbox}
\ex{1.19}{A set is infinite iff it contains a copy of $\N$}
\textbf{($\Leftarrow$)} If $f:\N\to S$ is injective, its outputs $f(1),f(2),\ldots$ are all distinct. A finite $S$ cannot contain infinitely many different elements (pigeonhole principle), so $S$ is infinite.

\textbf{($\Rightarrow$)} Suppose $S$ is infinite. Choose $s_1\in S$. After selecting distinct $s_1,\ldots,s_n$, the complement $S\setminus\{s_1,\ldots,s_n\}$ is nonempty: otherwise $S$ would be finite. Choose $s_{n+1}$ from this complement. Recursively obtain distinct elements and define $f(n)=s_n$. Then $f$ is injective. \textbf{Set-theoretic nuance:} the unrestricted infinite-set-to-injection direction uses a choice principle (typically dependent choice) in this recursive construction. In standard mathematics with AC, as in these lecture notes, the argument is valid; without choice, infinite sets need not contain a countably infinite subset. $\square$

\begin{tipbox}{Suggested practice total: mark-by-mark checklist}
\begin{tabular}{@{}p{0.77\linewidth}r@{}}
Injection implies infinite & 3 marks \\
Build countably many distinct elements using choice & 5 marks \\
Explain injectivity and conclusion & 2 marks \\
\hline\textbf{Total} & \textbf{10 marks}\\
\end{tabular}
\end{tipbox}
\newpage\section{Quick-reference formula sheet and proof templates}
\begin{keybox}{Everything worth memorising at a glance}
\begin{align*}
&d(x,y)=0\iff x=y,\quad d(x,y)=d(y,x),\quad d(x,z)\le d(x,y)+d(y,z);\\
&d_{\rm disc}(x,y)=\mathbf1_{x\ne y},\quad d_{\rm count}(A,B)=|A\triangle B|;\\
&d_1(f,g)=\int_0^1|f-g|,\quad d_\infty(f,g)=\sup_{t\in[0,1]}|f(t)-g(t)|;\\
&d_H(K,J)=\inf\{\delta\ge0:K^\delta\supseteq J,\ J^\delta\supseteq K\};\\
&A\triangle B=(A\setminus B)\cup(B\setminus A);\\
&f\text{ injective}\iff[f(x)=f(y)\Rightarrow x=y];\\
&f\text{ surjective}\iff (\forall t\in T)(\exists s\in S)\ f(s)=t;\\
&f^{-1}[C]=\{s:f(s)\in C\};\\
&f[\bigcup A_i]=\bigcup f[A_i],\quad f[\bigcap A_i]\subseteq\bigcap f[A_i];\\
&f^{-1}[\bigcup C_i]=\bigcup f^{-1}[C_i],\quad f^{-1}[\bigcap C_i]=\bigcap f^{-1}[C_i];\\
&f^{-1}[T\setminus C]=S\setminus f^{-1}[C],\quad |\mathcal F(S;T)|=|T|^{|S|}\ \text{(finite)}.
\end{align*}
\end{keybox}
\subsection{Which method should I use?}
\textbf{``Show this is a metric.''} Verify MS0, MS1a, MS1b, MS2 and MS3 separately; use cases, absolute values, cardinality inclusion, a supremum or integral as appropriate.

\textbf{``Show two sets are equal.''} Start ``Take $x$ in the left-hand side'' and show it belongs to the right; reverse direction. Or build a chain of equivalent membership statements.

\textbf{``Find the range / inverse.''} Set $y=f(x)$, solve for $x$ and identify forbidden outputs. If a left inverse must be defined outside the range, supply arbitrary legal values there.

\textbf{``Prove countability.''} Exhibit a list/bijection, inject into a known countable set, or express the set as a countable union of countable sets.

\textbf{``Show uncountability.''} Use Cantor's theorem and an injection from a known uncountable set; alternatively construct a diagonal contradiction.

\textbf{``Prove equivalence relation.''} Check reflexive, symmetric and transitive \emph{one by one}; if asked whether a condition always holds, one valid counterexample is enough to disprove it.
\section*{Source and scope note}
Both supplied lecture PDFs are by MATH6310, Part I, Sections 0 and 1, marked ``Last edited: 07.ix.2026''. This study aid retains the notes' conventions and sequence of topics. It expands explanations and supplies original full arguments for the numbered exercises. Optional results are marked as such; the statement of Schr\"oder--Bernstein is specifically required by the lecture notes, while its proof is labelled non-examinable there. The source PDFs themselves remain the authority for course-specific assessment guidance.
\end{document}
